Separable network games can capture situations in which players distribute their capacities or resources in a continuous way across multiple targets to achieve control over them. The targets can be battlefields, inspection sites, exit points etc.

In the class of games considered in this chapter, each player $i \in N$ is identified with a node of an undirected graph $(N,E)$ without loops. The presence of an edge $\{i,j\}\in E$ means that players $i$ and $j$ engage in a two-person general-sum strategic game in which the utility functions of players $i$ and $j$ are functions $u_{ij}\colon X_i\times X_j\to\R$ and $u_{ji}\colon X_j\times X_i\to\R$.

We will always assume that the functions $u_{ij}$ and $u_{ji}$ are at least integrable, ensuring that every expected utility appearing below is correctly defined. Typically, a player $i$ participates in multiple pairwise games, in which case, the same strategy strategy $x_i\in X_i$ is implemented in all two-person games played with the adjacent players in the graph $(N,E)$. Thus, the aggregated utility of player $i$ is computed as
\begin{equation}\label{def:pi}
    u_i(\bx) = \sum_{\substack{j\in N \\ \{i,j\}\in E}} u_{ij}(x_i,x_j), \qquad \bx\in\bX.
\end{equation}
The resulting $n$-person game with compact strategy sets is called a \emph{separable network game}, or ``network game'' for short. We say that $\calG$ is
\begin{itemize}
    \item a \emph{polymatrix game} when all strategy sets~$X_i$ are finite;
    \item \emph{zero-sum} if  $\sum_{i\in N} u_i(\bx)=0$ for all $\bx\in\bX$;
    \item \emph{continuous} if each utility function $u_i$ is continuous.
\end{itemize}
A class of network games often studied in the literature is the class of zero-sum polymatrix games. The ``global'' zero-sum assumption and finite strategy sets make such games computationally tractable \cite{Cai16,FPNet,Bailey19,Rubinstein,CaiDaskalkis}.

\begin{example}
Let $N=\{1,2,3,4\}$ and the graph $(N,E)$ be the cycle on \cref{fig:4}. This is the smallest example of a network game such that each player interacts only with a strict subset of other players. For example, the utility function of player $3$ is $u_3(x_1,x_2,x_3,x_4)=u_{31}(x_3,x_1)+u_{34}(x_3,x_4)$. The zero-sum condition for this game can be formulated as $u_3=-u_1-u_2-u_4$.
\end{example}

\begin{figure}[ht]
    \centering
  \begin{tikzpicture}
    \node[circle,fill=blue!20,draw,minimum size=0.5cm,inner sep=0pt] (1) {$1$};
    \node[circle,fill=blue!20,draw,minimum size=0.5cm,inner sep=0pt] (2) [right = 1cm  of 1]  {$2$};
    \node[circle,fill=blue!20,draw,minimum size=0.5cm,inner sep=0pt] (3) [below = 1cm  of 1] {$3$};
    \node[circle,fill=blue!20,draw,minimum size=0.5cm,inner sep=0pt] (4) [right = 1cm  of 3] {$4$};

    \path[draw,thick]
    (1) edge node {} (2)
    (1) edge node {} (3)
    (4) edge node {} (2)
    (3) edge node {} (4)
    ;
\end{tikzpicture}
    \caption{A network game with $4$ players and $4$ pairwise games}
    \label{fig:4}
\end{figure}

As a direct consequence of Glicksberg's theorem \ref{thm:Glick}, every continuous network game $\calG$ has a NE $\bmu^\star \in \bM$ in mixed strategies. We will show that the NE in a zero-sum continuous network game can be characterized using an optimization-based approach. Specifically, we consider the following infinite-dimensional linear optimization problem, which generalizes \emph{Linear Program 1} of \textcite{Cai16}:
\begin{mini!}{\bmu,\bw}{\sum_{i\in N} w_i} {\label{P}}{}
    \addConstraint{w_i}{\geq U_i(x_i,\bmu_{-i})\qquad \forall i\in N\enskip \forall x_i\in X_i} \label{P:ineq}
    \addConstraint{\bmu}{=(\mu_1,\dots,\mu_n) \in\bM}
    \addConstraint{\bw}{=(w_1,\dots,w_n) \in\R^n}
\end{mini!}

\begin{theorem}\label{thm:opti}
    For every profile of mixed strategies $\bmu^\star \in \bcalP$ in a zero-sum continuous network game $\calG$, the following are equivalent:
    \begin{enumerate}
        \item $\bmu^\star$ is a Nash equilibrium of game $\calG$.
        \item $(\bmu^\star, \bw^\star)$ is an optimal solution to (\ref{P}), where
        \begin{equation}\label{def:wi}
            w_i^\star = \max_{x_i \in X_i} U_i(x_i, \bmu_{-i}^\star), \qquad i \in N.
        \end{equation}
    \end{enumerate}
\end{theorem}

\begin{proof}
    Firstly, observe that the expression $\max_{\nu_i \in P_i} U_i(\nu_i, \bmu_{-i})$ is maximized by a $\nu_i$ supported on any of $U_i$'s maximizers. Therefore, given a feasible $(\bmu, \bw)$, the constraint (\ref{P:ineq}) implies
    \begin{equation}\label{ineq1}
        w_i \geq \max_{x_i \in X_i} U_i(x_i, \bmu_{-i}) = \max_{\nu_i \in P_i} U_i(\nu_i, \bmu_{-i}) \geq U_i(\mu_i, \bmu_{-i}) \quad \forall i \in N.
    \end{equation}
    Note that the maxima in inequalities (\ref{ineq1}) exist by compactness of $X_i$ and the continuity of the functions $U_i(\cdot, \bmu_{-i}) \colon X_i \to \R$. Applying inequalities (\ref{ineq1}) in the objective, together with the zero-sum property, gives:
    \begin{equation}\label{ineq3}
        \sum_{i \in N} w_i \geq \sum_{i \in N} U_i(\mu_i, \bmu_{-i}) = 0.
    \end{equation}

    By \cref{thm:Glick}, there exists an NE $\bmu^\star$ such that $U_i(\bmu^\star) = \max_{x_i \in X_i} U_i(x_i, \bmu_{-i}^\star) \; \forall i \in N$. Let $\bmu$ be an NE and $\bw^\star_i$ be as in condition (\ref{def:wi}), then the inequalities in (\ref{ineq1}) become equalities, and the lower bound of zero is achieved in (\ref{ineq3}).

    We have shown that NE correspond to optimal solutions. Conversely, consider any optimal solution $(\bmu^\star, \bw^\star)$. Notice that $w_i^\star$ must be tight upper bounds on $U_i(\cdot, \bmu_{-i}^\star)$. By the zero-sum property, $\sum_{i \in N} U_i(\bmu) = 0$ for any $\bmu$. If $\bmu^\star$ were not an NE, then there exists a player $i$ such that $\max_{x_i \in X_i} U_i(x_i, \bmu_{-i}) > U_i(\bmu)$. This implies that
    \[
        0 = \sum_{i \in N} U_i(\bmu) < \sum_{i \in N} \max_{x_i \in X_i} U_i(x_i, \bmu_{-i}) = \sum_{i \in N} w_i,
    \]
    i.e., the objective value is not optimal.
\end{proof}

The Borel probability measures and continuous functions in Problem~\ref{P} do not generally allow for effective computation. Still, there are cases where the solution can be computed. For example, in the case of zero-sum polymatrix games, Problem~\ref{P} becomes a linear program~\cite{Cai16}. Similarly, if we restrict the model to polynomial functions and semialgebraic sets, the solution can be computed using the tools of polynomial optimization discussed in \cref{sec:momsos}.